
\chapter{Introduction générale}
\label{ch:introduction}
\section{Contexte général}

À l'aube de 2026, les agglomérations urbaines connaissent
une croissance démographique sans précédent. L'augmentation
de la population entraîne mécaniquement une multiplication
du parc automobile, faisant de la congestion routière l'un
des défis les plus pressants des métropoles modernes.
Les heures de pointe paralysent quotidiennement des milliers
de citoyens, génèrent une pollution atmosphérique croissante
et compromettent l'intervention rapide des services d'urgence.

Face à ces enjeux, le concept de \textbf{ville intelligente}
(\textit{smart city}) s'impose comme une réponse structurée,
visant à mobiliser les technologies numériques au service de
la gestion urbaine. Le Maroc s'inscrit pleinement dans cette
dynamique : la ville d'Agadir a lancé en 2022 son programme
officiel de transformation en ville intelligente, initié
dans le cadre du programme de développement urbain impulsé
par Sa Majesté le Roi Mohammed VI. Cette initiative
témoigne d'une volonté politique forte de moderniser
les infrastructures urbaines marocaines~\cite{elalaoui2025}.

La gestion du trafic constitue l'une des problématiques
centrales de cette transformation. Dans les intersections
d'Agadir, comme dans de nombreuses villes marocaines,
les systèmes de feux à cycles fixes demeurent la norme.
Ces dispositifs, incapables de s'adapter aux conditions
réelles du trafic, engendrent des temps d'attente inutiles,
des embouteillages récurrents aux heures de pointe et
des situations dangereuses pour les véhicules d'urgence.
La présence récurrente d'agents de police aux carrefours
pour assurer une régulation manuelle illustre les limites
criantes des systèmes actuels.

L'essor de l'Intelligence Artificielle ouvre de nouvelles
perspectives pour résoudre ce problème. Des systèmes
capables de percevoir visuellement l'état d'une intersection,
d'analyser la densité du trafic en temps réel et de prendre
des décisions adaptatives constituent désormais une réalité
technologique accessible. Cependant, le déploiement de
caméras de surveillance dans l'espace public soulève une
question fondamentale : \textbf{comment garantir que ces
solutions respectent le droit à la vie privée des citoyens,
conformément à la Loi 09-08 relative à la protection
des personnes physiques à l'égard du traitement
des données à caractère personnel~?}

C'est précisément à cette double problématique —
efficacité de la gestion du trafic et respect de la
vie privée — que ce mémoire apporte une réponse
technique et scientifique.

\section{Motivation}

\subsection{Pourquoi l'Edge AI ?}

Les solutions traditionnelles de surveillance du trafic
reposent généralement sur une architecture Cloud :
les flux vidéo des caméras sont transmis vers des
serveurs distants pour traitement, puis les décisions
sont renvoyées vers les équipements locaux. Cette
approche présente plusieurs inconvénients majeurs
dans notre contexte.

D'une part, la transmission continue de flux vidéo
vers le Cloud engendre une latence incompatible avec
les exigences du temps réel. D'autre part, et c'est
le point le plus critique, \textbf{chaque image transmise
constitue une donnée personnelle} : elle contient
des visages, des plaques d'immatriculation et des
informations comportementales sur les citoyens.
Transmettre ces données vers des serveurs distants
expose à des risques majeurs de violation de la
vie privée et contredit les principes fondamentaux
de la Loi 09-08.

L'\textbf{Edge AI} apporte une réponse structurelle
à ces deux problèmes : le traitement s'effectue
\textit{localement}, sur le dispositif même, sans
qu'aucune donnée personnelle ne soit transmise
vers l'extérieur. Les données restent là où elles
sont collectées.

\subsection{Pourquoi le Raspberry Pi 5 ?}

Le choix d'un dispositif embarqué à ressources
limitées — le Raspberry Pi 5 — n'est pas une
contrainte subie mais un choix délibéré et justifié.
Il répond à trois impératifs concrets :

\begin{itemize}
    \item \textbf{Accessibilité économique :}
    un déploiement à l'échelle d'une ville
    nécessite un équipement pour chaque
    intersection. Le coût unitaire d'un
    Raspberry Pi 5 est sans commune mesure
    avec celui d'un serveur dédié.

    \item \textbf{Déployabilité :}
    sa compacité et sa faible consommation
    électrique le rendent directement
    intégrable dans les armoires techniques
    existantes des carrefours.

    \item \textbf{Contribution scientifique :}
    à notre connaissance, aucun travail existant
    ne documente le déploiement de YOLO26n
    sur Raspberry Pi 5 dans un contexte
    de gestion du trafic urbain.
\end{itemize}

\subsection{Pourquoi le Privacy-by-Design ?}

L'évolution rapide des technologies de vision
par ordinateur et la multiplication des dispositifs
de surveillance ont rendu la protection des données
personnelles plus urgente que jamais. Des travaux
récents ont démontré qu'il est possible de
ré-identifier des individus même depuis des images
partiellement anonymisées~\cite{razi2025,jeremiah2024}.

Face à ce constat, le principe de
\textbf{\textit{Privacy-by-Design}} s'impose
comme la seule approche véritablement protectrice :
la protection des données n'est pas ajoutée
\textit{a posteriori} comme une fonctionnalité
optionnelle, mais intégrée dès la conception
au cœur de l'architecture du système.

Dans notre pipeline, cela se traduit concrètement
par l'application du floutage gaussien
\textbf{immédiatement après la détection},
avant tout traitement ultérieur, de sorte
qu'aucune donnée personnelle identifiable
ne circule dans le système — conformément
aux exigences de la Loi 09-08 marocaine.

\section{Problématique}

La littérature scientifique récente témoigne d'efforts
significatifs dans trois directions complémentaires.
En matière de détection d'objets, les approches
\textit{one-stage} — dont la famille YOLO constitue
le représentant le plus emblématique — ont démontré
leur supériorité en termes de vitesse d'inférence,
tandis que les approches \textit{two-stage} offrent
une précision accrue au détriment de la latence
\cite{ramos2025}. Sur le plan du déploiement embarqué,
les travaux récents explorent des variantes nano
et des techniques de quantification pour adapter
ces modèles aux dispositifs à ressources limitées
\cite{wang2025, sapkota2025}. Enfin, concernant
la protection des données, des techniques
d'anonymisation par floutage ou remplacement
de silhouette ont été proposées pour les systèmes
de vidéosurveillance \cite{jeremiah2024, razi2025}.

Cependant, une analyse approfondie de l'état de l'art
révèle une lacune fondamentale : \textbf{aucun travail
existant ne propose un pipeline complet et intégré}
combinant simultanément détection temps réel,
anonymisation Privacy-by-Design et décision adaptative
des feux, déployé sur un dispositif embarqué
minimaliste, et validé sur des données locales
représentatives du contexte urbain marocain.
Michailidis et al. confirment d'ailleurs que
\textit{"le déploiement en conditions réelles demeure
le défi le plus pressant et non résolu"} dans
ce domaine~\cite{michailidis2025}.

\begin{definitionbox}{Question centrale de recherche}
Comment concevoir et déployer un système intelligent
de gestion adaptative des feux de circulation en
temps réel, sur dispositif embarqué à ressources
limitées, tout en garantissant la protection des
données personnelles des usagers, dans le contexte
des smart cities marocaines~?
\end{definitionbox}

\section{Objectifs de recherche}

Pour répondre à la problématique posée, ce travail
poursuit cinq objectifs complémentaires et interdépendants :

\begin{tcolorbox}[
  colback  = PrimaryBlue!5,
  colframe = PrimaryBlue,
  arc      = 4pt,
  boxrule  = 1pt
]

\textbf{O1 — Détection multi-classes en temps réel}\\
Développer et fine-tuner un modèle de détection d'objets
capable d'identifier en temps réel les véhicules
(voitures, bus, camions, motos, véhicules d'urgence)
et les piétons dans un contexte de surveillance
urbaine marocain.

\vspace{0.4cm}

\textbf{O2 — Décision adaptative des feux}\\
Concevoir un module de décision adaptative des phases
de feux de circulation, basé sur la densité de trafic
détectée et formalisé comme un Processus de Décision
Markovien à politique déterministe.

\vspace{0.4cm}

\textbf{O3 — Déploiement embarqué temps réel}\\
Déployer le pipeline complet de perception et de
décision en temps réel sur un dispositif embarqué
à ressources limitées (Raspberry Pi 5), via
export ONNX, seuils d'inférence adaptés et
optimisations CPU. La validation temps réel porte
sur la famille YOLO26n nano, tandis que le modèle
Step 3 960 est retenu comme meilleur compromis de
précision.

\vspace{0.4cm}

\textbf{O4 — Conformité Loi 09-08}\\
Garantir la conformité à la Loi 09-08 relative
à la protection des données personnelles, via
l'anonymisation automatique des piétons détectés
selon le principe Privacy-by-Design.

\vspace{0.4cm}

\textbf{O5 — Validation sur données réelles}\\
Valider les performances du système sur des données
réelles collectées dans le contexte urbain marocain,
notamment à Lqliaa (Agadir).

\end{tcolorbox}

\section{Hypothèses de recherche}

La démarche scientifique adoptée dans ce mémoire
repose sur cinq hypothèses de recherche vérifiables
empiriquement. Chacune est associée à un objectif
et sera validée ou infirmée par les résultats
expérimentaux présentés au Chapitre~\ref{ch:resultats}.

\begin{warningbox}{H1 — Détection en temps réel}
Nous supposons qu'un modèle de détection d'objets
de type \textit{one-stage}, fine-tuné sur un dataset
incluant des données locales marocaines, peut atteindre
des performances suffisantes pour la détection de
véhicules et piétons en conditions réelles urbaines.
\end{warningbox}

\vspace{0.3cm}

\begin{warningbox}{H2 — Apport du dataset hybride}
Nous supposons qu'un fine-tuning intégrant des données
locales marocaines, CityPersons et des crops de petits
piétons améliore la robustesse du modèle dans le
contexte routier étudié.
\end{warningbox}

\vspace{0.3cm}

\begin{warningbox}{H3 — Déploiement embarqué temps réel}
Nous supposons qu'un dispositif embarqué à ressources
limitées, via l'optimisation du modèle de détection,
peut atteindre une précision suffisante pour la gestion
du trafic et une latence compatible avec le traitement
en temps réel.
\end{warningbox}

\vspace{0.3cm}

\begin{warningbox}{H4 — Protection des données}
Nous supposons que l'application d'une technique
de protection des données personnelles en temps réel
sur les piétons détectés n'impacte pas significativement
les performances globales du système.
\end{warningbox}

\vspace{0.3cm}

\begin{warningbox}{H5 — Décision adaptative}
Nous supposons qu'une fonction de décision basée
sur l'état courant de l'intersection — défini par
la densité et le type des objets détectés — permet
de générer automatiquement une politique d'allocation
des phases de feux qui réduit le temps d'attente moyen
et garantit la priorité aux véhicules d'urgence.
\end{warningbox}

\vspace{0.5cm}

\begin{table}[H]
\centering
\caption{Correspondance hypothèses — objectifs — métriques}
\begin{tabularx}{\textwidth}{clXl}
\toprule
\textbf{H} & \textbf{Hypothèse} &
\textbf{Métrique de validation} &
\textbf{Objectif} \\
\midrule
H1 & Détection temps réel &
mAP@50 test $\geq$ 75\% & O1 \\
H2 & Dataset hybride &
Gain Step 3 800 vs 960 et couverture person & O1 \\
H3 & Edge embarqué &
YOLO26n nano $\geq$ 5 FPS; Step 3 800 = 3.45--3.94 FPS & O3 \\
H4 & Protection données &
100\% des détections person floutées & O4 \\
H5 & Décision adaptative &
Validation 5 scénarios & O2 \\
\bottomrule
\end{tabularx}
\label{tab:hypotheses}
\end{table}
\section{Contributions originales}

Ce travail apporte six contributions originales
à l'état de l'art, documentées et vérifiables :

\begin{resultbox}{C1 — Premier benchmark YOLO26n sur Raspberry Pi 5}
À notre connaissance, ce travail constitue le premier
benchmark empirique de YOLO26n — architecture publiée
en septembre 2025 — sur Raspberry Pi 5 dans un contexte
de gestion du trafic urbain, comparant les formats
FP32 et INT8 sur CPU ARM Cortex-A76.
\end{resultbox}

\vspace{0.3cm}

\begin{resultbox}{C2 — Optimisation précision-vitesse simultanée}
Contrairement aux travaux existants qui privilégient
soit la précision (architectures lourdes sur GPU),
soit la vitesse (modèles très allégés avec perte
de mAP), notre approche combine un modèle accuracy
\textbf{YOLO26n Step 3 960} atteignant
\textbf{79,2\% mAP50 test}; pour le déploiement
embarqué, la famille YOLO26n nano est validée à
\textbf{7,6 FPS réels} sur Raspberry Pi 5.
\end{resultbox}

\vspace{0.3cm}

\begin{resultbox}{C3 — Dataset hybride avec données locales marocaines}
Face à l'absence de datasets publics annotés
représentatifs du contexte urbain marocain,
ce travail constitue et documente un dataset
hybride de \textbf{6 761 images} combinant
sources internationales et données locales
collectées à \textbf{Lqliaa, Agadir}.
Le dataset final contient notamment
\textbf{23 096 instances person}, utilisées pour
renforcer la couverture d'anonymisation.
\end{resultbox}

\vspace{0.3cm}

\begin{resultbox}{C4 — Pipeline complet embarqué sur dispositif unique}
Ce travail propose un pipeline complet
intégrant perception, anonymisation et décision
adaptative sur un \textbf{dispositif unique}
(Raspberry Pi 5), sans recours au Cloud,
validant ainsi la faisabilité d'une architecture
entièrement embarquée pour la gestion intelligente
du trafic urbain.
\end{resultbox}

\vspace{0.3cm}

\begin{resultbox}{C5 — Décision adaptative déterministe embarquée}
Le module de décision MDP à politique déterministe
proposé offre une alternative computationnellement
légère aux approches par Apprentissage par
Renforcement, compatible avec les contraintes
CPU ARM du Raspberry Pi 5, tout en garantissant
la priorité absolue aux véhicules d'urgence.
\end{resultbox}

\vspace{0.3cm}

\begin{resultbox}{C6 — Privacy-by-Design sur Edge Device}
L'anonymisation des piétons est appliquée
\textbf{immédiatement après la détection},
avant tout traitement ultérieur, selon le principe
Privacy-by-Design. Le système floute
\textbf{100\% des boîtes détectées} de classe
\textit{person} avant tracking. La limite restante
concerne les personnes non détectées par le modèle,
ce qui est explicitement documenté dans l'évaluation.
\end{resultbox}

\vspace{0.5cm}

\begin{table}[H]
\centering
\caption{Synthèse des contributions originales}
\begin{tabularx}{\textwidth}{clX}
\toprule
\textbf{C} & \textbf{Contribution} &
\textbf{Élément distinctif} \\
\midrule
C1 & Benchmark YOLO26n Pi 5 &
Premier test empirique sur ce hardware \\
C2 & Précision + vitesse &
0,792 mAP50 test; YOLO26n nano à 7,6 FPS Pi 5 \\
C3 & Dataset marocain &
6 761 images, 23 096 instances person \\
C4 & Pipeline complet embarqué &
Perception + anonymisation + décision \\
C5 & Décision MDP déterministe &
Légère, temps réel, priorité urgence \\
C6 & Privacy-by-Design edge &
100\% des détections person floutées \\
\bottomrule
\end{tabularx}
\label{tab:contributions}
\end{table}
\newpage
\section{Planification et méthodologie}

\subsection{Méthodologie de travail}

Ce projet a été conduit selon une méthodologie
structurée s'appuyant sur des pratiques de
traçabilité expérimentale et sur une couche
\textbf{DataOps/MLOps légère}. L'objectif n'est
pas de prétendre à une infrastructure industrielle
complète, mais de rendre explicites les versions
du dataset, les paramètres d'entraînement, les
métriques finales et les artefacts de déploiement.
Cette traçabilité est matérialisée par
\texttt{dvc.yaml}, \texttt{dvc.lock},
\texttt{params.yaml},
\texttt{metrics/final\_yolo26n\_step3\_960.json},
un artefact dataset suivi par DVC, et une
journalisation MLflow locale adossée à
\path{mlflow.db}.

\begin{table}[H]
\centering
\caption{Outils utilisés pour la traçabilité expérimentale}
\begin{tabularx}{\textwidth}{llX}
\toprule
\textbf{Outil} & \textbf{Usage} &
\textbf{Justification} \\
\midrule
Git/GitHub & Versioning code et gestion de projet &
Branches \texttt{main}/\texttt{develop}, historique des commits
et préparation des livrables stables \\
Pipeline DVC & DataOps dataset &
Suivi du dataset final et formalisation des étapes :
split sans fuite, ajout CityPersons et crops tiny persons \\
MLflow local & MLOps expérimental &
Journalisation SQLite des paramètres, métriques
et artefacts finaux \\
YAML/JSON & Configuration et métriques &
Paramètres reproductibles et résultats finaux
lisibles par humains et scripts \\
Google Colab Pro & Fine-tuning &
GPU NVIDIA A100-SXM4-40GB \\
Ultralytics & Entraînement et validation &
Courbes, métriques, poids \texttt{best.pt} \\
Roboflow & Annotation &
Dataset Lqliaa et sources externes \\
ONNX Runtime & Déploiement &
Inférence CPU sur Raspberry Pi 5 \\
OpenCV + Flask & Pipeline et dashboard &
Traitement vidéo, anonymisation et supervision \\
\bottomrule
\end{tabularx}
\label{tab:mlops}
\end{table}

Dans cette organisation, la partie \textbf{DataOps}
correspond à la traçabilité des transformations du
dataset, la partie \textbf{MLOps} au suivi des
expériences et des métriques, et la partie
\textbf{DevOps} au passage reproductible du modèle
entraîné vers le pipeline embarqué ONNX + Flask.
Le code source et le rapport sont structurés dans
un dépôt GitHub avec une branche \texttt{main}
réservée aux versions stables et une branche
\texttt{develop} utilisée pour l'intégration des
corrections, scripts, figures et évolutions du
pipeline.
Les données volumineuses restent hors Git ; le
dataset final est suivi par DVC via
\path{data/dataset/dataset_step3_tiny_person_crops.dvc},
et l'expérience finale est journalisée dans MLflow
via \path{mlflow.db}.

La démarche scientifique adoptée repose sur
un processus itératif en cinq étapes :

\begin{enumerate}
    \item \textbf{Recherche bibliographique :}
    sélection et analyse de 12 articles
    scientifiques Q1-Q4 couvrant les six
    piliers thématiques du projet.

    \item \textbf{Expérimentation contrôlée :}
    benchmark empirique de trois familles
    YOLO (v8, v11 et v26), en variantes nano
    et small, sur GPU puis sur Raspberry Pi 5.

    \item \textbf{Validation par hypothèses :}
    chaque décision technique est ancrée
    dans une hypothèse vérifiable empiriquement.

    \item \textbf{Traçabilité DataOps/MLOps :}
    formalisation des étapes dataset dans
    \texttt{dvc.yaml}, centralisation des
    hyperparamètres dans \texttt{params.yaml}
    et conservation des métriques finales au
    format JSON.

    \item \textbf{Déploiement progressif :}
    développement sur Mac, export ONNX,
    puis déploiement sur Raspberry Pi 5.
\end{enumerate}
\newpage
\subsection{Planification du projet}

Le projet s'est déroulé de \textbf{mars 2026}
à \textbf{juin 2026}, organisé en huit phases
séquentielles illustrées par le diagramme
de Gantt en Figure~\ref{fig:gantt}.

\begin{figure}[H]
\centering
\resizebox{0.96\textwidth}{!}{%
\begin{tikzpicture}[
  font=\small,
  month/.style={
    draw=DarkGray!45, fill=LightGray,
    minimum width=2.25cm, minimum height=0.65cm,
    align=center, font=\bfseries\small
  },
  task/.style={
    rounded corners=4pt,
    minimum height=0.46cm,
    align=center,
    font=\scriptsize\bfseries,
    text=white
  },
  label/.style={font=\scriptsize, align=right, text width=3.9cm},
  gridline/.style={DarkGray!25, line width=0.35pt}
]

% Axe temporel
\foreach \x/\m in {0/Mars,2.25/Avril,4.5/Mai,6.75/Juin} {
  \node[month] at (\x,0) {\m};
  \draw[gridline] (\x-1.125,-0.40) -- (\x-1.125,-5.85);
}
\draw[gridline] (7.875,-0.40) -- (7.875,-5.85);

% Lignes des tâches
\foreach \y in {-0.95,-1.55,-2.15,-2.75,-3.35,-3.95,-4.55,-5.15} {
  \draw[gridline] (-6.25,\y-0.28) -- (7.9,\y-0.28);
}

\node[label] at (-4.15,-0.95) {État de l'art};
\node[task, fill=PrimaryBlue, minimum width=2.95cm] at (-0.25,-0.95) {};

\node[label] at (-4.15,-1.55) {Conception M1--M6};
\node[task, fill=PrimaryBlue!82, minimum width=2.55cm] at (1.00,-1.55) {};

\node[label] at (-4.15,-2.15) {Dataset et annotation};
\node[task, fill=AccentGreen!82!black, minimum width=3.45cm] at (2.00,-2.15) {};

\node[label] at (-4.15,-2.75) {Fine-tuning YOLO26n};
\node[task, fill=AccentGreen!70!black, minimum width=2.85cm] at (3.25,-2.75) {};

\node[label] at (-4.15,-3.35) {Pipeline embarqué};
\node[task, fill=AccentOrange, minimum width=3.10cm] at (4.40,-3.35) {};

\node[label] at (-4.15,-3.95) {Dashboard et anonymisation};
\node[task, fill=AccentOrange!88!black, minimum width=2.55cm] at (5.20,-3.95) {};

\node[label] at (-4.15,-4.55) {Benchmark Pi 5};
\node[task, fill=red!65, minimum width=1.85cm] at (6.35,-4.55) {};

\node[label] at (-4.15,-5.15) {Rédaction et validation};
\node[task, fill=DarkGray, minimum width=3.05cm] at (6.10,-5.15) {};

\node[
  draw=DarkGray!45,
  fill=white,
  rounded corners=4pt,
  font=\scriptsize,
  align=center,
  text width=8.9cm
] at (2.15,-6.10)
{Lecture : les tâches se chevauchent volontairement, car les résultats
du benchmark et du dashboard ont alimenté directement la rédaction finale.};
\end{tikzpicture}}
\caption{Diagramme de Gantt — Planification du projet}
\label{fig:gantt}
\end{figure}
\newpage
% ── Section 1.8 ──────────────────────────────────────────────
\section{Organisation du mémoire}

Ce mémoire est structuré en sept chapitres
complémentaires, dont l'enchaînement reflète
la démarche scientifique adoptée :

\textbf{Le Chapitre 2} présente l'état de l'art
à travers six piliers thématiques. Il expose
les fondements théoriques de la détection
d'objets (approches one-stage et two-stage,
évolution de YOLO v1 à v26), de l'Edge AI
et de l'optimisation embarquée, des techniques
de protection des données personnelles,
et de la décision adaptative par processus
de Markov. Il se conclut par l'identification
du gap scientifique que ce travail comble.

\textbf{Le Chapitre 3} présente l'analyse
des besoins et la conception du système.
Il formalise l'architecture du pipeline M1--M6,
le modèle MDP de décision adaptative avec
ses formules mathématiques, et justifie
chaque choix technologique par rapport
à l'état de l'art.

\textbf{Le Chapitre 4} détaille l'implémentation
complète du système : construction du dataset
hybride, fine-tuning des trois architectures
YOLO, benchmarks sur GPU et Raspberry Pi 5,
et développement de chaque module du pipeline.

\textbf{Le Chapitre 5} présente les résultats
expérimentaux et valide les cinq hypothèses
de recherche à travers des métriques précises,
démontrant le fonctionnement du pipeline
complet sur données réelles marocaines.

\textbf{Le Chapitre 6} discute les contributions,
les limites du système et les perspectives
de recherche futures.

\textbf{Le Chapitre 7} présente la conclusion
générale du mémoire, synthétise les résultats
obtenus et répond explicitement à la problématique
posée.
